\subsection{Les outils du dépôts légal : du sur mesure pour la nouvelle mission patrimoniale}

%% Introduction %%

À l'arrivée du dépôt légal, l'INA possède donc déjà un outil pour la gestion des données : IMAGO, mais il est estimé dépassé et trop rigide pour s'adapter à la nouvelle logique du dépôt légal. Il est nécessaire de mettre en place un outil qui convient aux besoins de cette nouvelle mission. 

%% Intégration du logiciel du dépôt légal : Hyperbase %%

On favorise l'utilisation de nouvelles bases de données. Une pour la radio : DLradio, une pour la télévision hertzienne : DLTV, une pour les chaînes du câble, du satellite et de la TNT : DLSat, une pour les chaînes régionales : DLREG. Un nouvel outil de production est mis en place : DLexe et un nouvel outil de consultation de la donnée : Hyperbase. Chaque champ a son référentiel. L'interface a été réalisé en adéquation avec les besoins de l'équipe du dépôt légal. Le dépôt légal fonctionnant en flux et pas en stock comme le font les archives professionnelles, les outils sont de ce fait concentrés sur les chaînes, puis les tranches horaires. Le fonctionnement des outils de production et de consultation devait être élaboré en fonction de ces prérequis. Ils fonctionnent avec trois interfaces : 
\begin{itemize}
	\item Gestion des journées de programmes : Contiennent l'ensemble des émissions d'une tranche horaire. Se présentant sous forme de liste. Aujourd'hui, c'est une fonctionnalité gérée par les TGDM.
	\item Gestion des matériels : Cette fonctionnalité évoluera bien évidemment avec la numérisation mais à la création du dépôt légal, les chaînes déposent encore le dépôt légal en supports physiques. 
	\item Saisie des documents\footcite[]{Rodes1994}
\end{itemize}
La traitement documentaire, quant à lui, se recentre sur la structure de l'émission. Les \enquote{notices de collection}, décrivant une émission, une collection et pas seulement une diffusion prend plus d'importance. Dans la même lignée. Dans l'indexation on donne une grande place à la description de décors techniques. Notion qui n'intéressait les correspondants de chaînes télévisées. 

%% IMAGO III et Totem %%

Les archives professionnelles étant en stock, les outils de production et de consultation de données sont plus concentrés sur l'émission, l’œuvre, le thème. L'unité est à un niveau inférieur à l'émission. La description est concentrée sur le \enquote{sujet} traité par l'émission car c'est ce qui sera recherché. L'indexation décrit avant tout l'image. Ceci dit, il est important de considérer que les notices pour les archives professionnelles ne sont pas créées ex nihilo, elles sont transmises de la base de données du dépôt légal, le traitement documentaire est donc adapté à des fins commerciales. 

\subsection{La recherche au dépôt légal et dans les archives professionnelles}

Ainsi, ces deux outils de consultation ne sont pas construits pour être utilisés de la même façon. Hyberbase, l'outil de consultation à disposition des documentalistes de l'INA et des chercheurs de l'INAthèque est construit pour combiner des requêtes afin de construire un corpus. Dans le sens inverse, \enquote{Totem}, l'outil de consultation et de production qui a succédé à \enquote{IMAGO} est tourné sur la recherche d'images pour les professionnels. Il est déjà organisé en corpus, élaboré par les documentalistes. Il est conçu pour lancer une première recherche que l'on peut affiner par la suite\footnote{(Voir entretien d'une documentaliste sur les bases de données de l'INA~\ref{ann:transcriptions},\nameref{ann:transcriptions}), p.~\pageref{ann:transcriptions}.}. 

Les outils de production des archives professionnelles et du dépôt légal sont particulièrement représentatifs de la séparation du traitement documentaire qui a opposé des années durant les services des archives professionnelles et du dépôt légal. Par leur fonctionnement, ils illustrent l'opposition des politiques documentaires qui ont opposé ces deux services des années durant et il contraignent les utilisateurs des bases de données internes à l'INA à toujours mobiliser ces deux politiques documentaires, ce qui ne facilite pas la compréhension des collections dans leur ensemble, tant elles sont décrites et présentées au public sous différents angles. 

\subsection{Le lac de données}

Cela fait de nombreuses années que l'INA a conscience de ces déséquilibres entrainés par la séparation des services et des outils du dépôt légal et des archives professionnelles. C'est ainsi que né le \enquote{lac de données} en 2015. Il a pour objectif de centraliser l'ensemble des données de l'INA en une seule infrastructure. Cette fois, il n'est plus question d'élaborer les outils en fonction des usages mais en fonction des données\footcite[]{institutnationaldhistoiredelartinha2019a}. En conséquent, un nouveau modèle de données est mis en place dans lequel toutes les données sont misent en relation et tous les référentiels des différents services de l'INA sont repris et mis en commun. Les collections sont divisées en trois entités, elles-mêmes décrites par des métadonnées : le contenu (les images, le son, l'émission en elle-même), l'évènement (date de diffusion, de production etc\dots) et le matériel (le support, le fichier sur qui contient le contenu)\footcite[87]{bassaïstéguy2025}. Ainsi, il doit être plus facile de convenir aux usages des données de tous les services de l'INA. 

Le lac permettrait de résoudre beaucoup des défauts que nous avons constaté précédemment. On opère une migration des données à partir de 2017 vers les nouvelles bases de données et on opère dans le même temps un énorme travail sur la qualité des données. Les données en double sont dédoublonnées, les erreurs corrigées, les données enrichies. Mais aujourd'hui, après plusieurs migrations des métadonnées et référentiels c'est un traitement différentiel qui est mis en place et encore une fois, il marque une différence entre les données du dépôt légal et des archives professionnelles. Aujourd'hui les données de la base de données des archives professionnelles sont synchronisées avec le lac de données toutes les 30 minutes quand ce sont seulement les référentiels du dépôt légal qui sont synchronisés, une fois par jour\footnote{\enquote{Le modèle de données de l'INA}, présentation powerpoint conçue par Axel Roche-Dioré, mis à jour par Soline Doat, 6 novembre 2025.}.  

Malgré les avantages que représente le lac de données pour notre outil, il n'est pour le moment pas possible de le mobiliser sans faire impasse sur une partie des données récentes du dépôt légal, sur lequel repose cette cartographie du traitement documentaire. 